\subsection{三项典型工程难点的排查与解决过程}

除了前文第~\ref{sec:cleaning} 节详述的 AST-v2 评分器重写外，团队在推进过程中还攻克了三项颇具代表性的工程难题：

\Needspace{15\baselineskip}
\begin{wrapfigure}{r}{0.31\linewidth}
\vspace{-8pt}
\centering
\shot[0.92\linewidth]{\detokenize{figures/screenshots/live-mobile-s4.jpg}}
\caption{窄屏手机视口下第四屏自动解除整屏吸附并实现零水平溢出堆叠}
\label{fig:mobile-s4}
\vspace{-8pt}
\end{wrapfigure}

\paragraph{难点一：多轮代码修复反馈构造器的防泄漏硬隔离}在设计 \texttt{code\_repair\_v2} 补充实验时，我们格外警惕一个致命的实验漏洞：如果评测脚本把公开诊断测试（\texttt{public}）和隐藏强保留测试（\texttt{held-out}）放在同一个字典里传给下一轮提示词拼接函数，哪怕只是不小心泄漏了一个隐藏报错行号或期望输出，多轮修复的高通过率就会沦为“看答案抄作业”。为此，我们在 \texttt{experiments/code\_repair\_v2/src/protocol.py} 中从接口签名层做了硬隔离：\texttt{feedback\_message(arm, evaluation)} 只被允许读取 \texttt{evaluation["public"]} 分支下的白名单投影字段（如诊断组只投影异常类型、行号与候选栈，反例组只投影输入输出对），并在 \texttt{test\_repair\_v2\_protocol.py} 中专门构造塞有哨兵密钥的伪造 \texttt{held-out} 字典做反泄漏单元测试，验证任何反馈臂生成的提示词中绝对不包含半个字节的隐藏测试信息。同时，当某份候选代码在模块导入阶段即因语法或顶层异常崩溃（\texttt{load\_failed=true}）时，分析器将“阻断未测用例数（\texttt{blocked\_cases}）”与“实际逐例执行失败数（\texttt{executed\_failed\_cases}）”分列记录，主指标严格按程序份数统计，彻底理顺了分母逻辑。

\paragraph{难点二：桌面整屏对焦与移动端窄屏流式排版的零溢出兼容}主站桌面端采用 CSS \texttt{scroll-snap-type: y mandatory} 实现类似光学镜头逐档对焦的整屏切换体验。但在真机模拟测试时我们发现，一旦到了高度不足 \texttt{680px} 的矮屏笔记本或宽度仅 \texttt{390px} 的手机竖屏上，一屏之内根本塞不下双栏图版与配对卡片，强行开启整屏吸附会导致下半截图表永远滚不到底。我们通过多维媒体查询（\texttt{@media (max-width: 960px), (max-height: 700px)}）实施了平滑降级：窄矮屏自动解除强制吸附、将左右分栏改为单列纵向流式堆叠，并在 \texttt{zhiwei.js} 中编写 \texttt{fitText} 字符宽度自适应裁剪与 SVG \texttt{viewBox="0 0 W H"} 等比缩放（实测效果见图~\ref{fig:mobile-s4}），最终在从 $320\times 568$ 到 $2560\times 1440$ 的全部测试视口下均做到了 \texttt{documentElement.scrollWidth === innerWidth}（零像素横向溢出）。

\paragraph{难点三：历史旧口径与新补充实验在前端的透明并存}由于早期案例短报告基于原 248,586 条统一表撰写，而主站第四屏接入了更严谨的 24,600 条代码反馈补充批次，如何防止读者在对照阅读时产生困惑？我们在数据快照层保留了双轨元数据，在主站首屏、第四屏图版 VII 读图说明及第五屏材料卡片处均显式注明两套批次编号与样本替换算式，同时把早期完整 44 图探索台原封不动保留在 \texttt{legacy.html} 中，做到了新旧成果互不覆盖、各司其职。
